\documentclass[10pt]{article}
\usepackage[margin=1in]{geometry}
\usepackage[T1]{fontenc}
\usepackage{lmodern,microtype,amsmath,amssymb,amsthm}
\theoremstyle{definition}
\newtheorem{assumption}{Assumption}
\begin{document}
\section*{Parameter-Free Composite Convex Optimization}
We minimize $F=f+\psi$ over a convex set $Q$ in a real Hilbert space $E$, where $f:E\to\mathbb{R}$ is differentiable and convex and $\psi:E\to\mathbb{R}$ is convex on $Q$. The continuous gradient satisfies generalized smoothness with hidden constants. A minimizer $x^\star\in Q$ has bounded distance and objective gap from the initial point $x^0\in Q$, and bounded gradient norm. Each call returns an exact gradient or an exact proximal minimizer. Only $E$, $x^0$, and a positive scalar $h_0$ are public. For hidden accuracy $\varepsilon>0$, the output $\bar{x}\in Q$ must satisfy $F(\bar{x})-F(x^\star)\leq\varepsilon$ almost surely.

The optimization problem is
\[
\min_{x\in Q} F(x), \qquad F(x)=f(x)+\psi(x).
\]
The norm on $E$ is induced by its inner product. The setup parameters satisfy
\[
h_0,L_0,R,\Delta,\varepsilon>0,
\qquad L_1,G\geq 0.
\]
Here $L_0,L_1$ control generalized smoothness, $R$ and $\Delta$ bound initialization, and $G$ bounds the gradient at a minimizer. The admissible instances satisfy the following conditions.

\begin{assumption}[Convexity and generalized smoothness]
The function $f$ is convex and differentiable everywhere on $E$, and $\nabla f:E\to E$ is continuous. For every $x,y\in E$ and $M\in\mathbb{R}$,
\[
\left(\|\nabla f(x+t(y-x))\|\leq M
\quad\text{for every }t\in[0,1]\right)
\ \Longrightarrow\ 
\|\nabla f(x)-\nabla f(y)\|
\leq (L_0+L_1M)\|x-y\|.
\]
\end{assumption}

\begin{assumption}[Composite structure and proximal attainment]
The set $Q\subseteq E$ is convex, and $\psi$ is convex on $Q$. For every $a>0$ and $z,v\in E$, a specified proximal map $\operatorname{prox}$ satisfies
\begin{align*}
\operatorname{prox}(a,z,v)&\in Q,\\
\operatorname{prox}(a,z,v)&\in
\mathop{\operatorname{argmin}}_{u\in Q}
\left\{
\frac{\|u-z\|^2}{2a}
+\langle v,u-z\rangle+\psi(u)
\right\}.
\end{align*}
\end{assumption}

\begin{assumption}[Attainment and initialization]
The initial point belongs to $Q$. There exists $x^\star\in Q$ such that
\begin{align*}
F(x^\star)&\leq F(x) \qquad\text{for every }x\in Q,\\
\|x^0-x^\star\|&\leq R,\\
F(x^0)-F(x^\star)&\leq\Delta,\\
\|\nabla f(x^\star)\|&\leq G.
\end{align*}
\end{assumption}

\begin{assumption}[Exact oracle access]
At every call index and for every realization of oracle randomness, a gradient query at $x\in E$ returns $\nabla f(x)\in E$, and a proximal query $(a,z,v)\in\mathbb{R}\times E\times E$ returns $\operatorname{prox}(a,z,v)\in E$. Each query of either kind costs one oracle call. The proximal minimization guarantee applies whenever $a>0$. Any internal algorithm randomness is independent of oracle randomness.
\end{assumption}

Queries may be adaptive, and their points need not belong to $Q$. The algorithm is selected using only $(E,x^0,h_0)$; the instance functions and feasible set are accessed through the oracle. The parameters $L_0,L_1,R,\Delta,G,\varepsilon$ are not supplied to the algorithm. An externally specified total call budget $K\in\mathbb{N}$ may depend on the full parameter tuple.

After $K$ calls, the algorithm returns $\bar{x}\in E$. Writing $\mathbb{P}$ for the joint law of algorithm and oracle randomness, success requires a minimizer $x^\star\in Q$ such that
\[
\bar{x}\in Q
\quad\text{and}\quad
F(\bar{x})-F(x^\star)\leq\varepsilon
\qquad\mathbb{P}\text{-almost surely}.
\]
\end{document}
